%-----------------------------------
% Kapitel 6: Fazit
%-----------------------------------
\section{Fazit}
\label{fazit}

Abschnitt~\ref{zusammenfassung} verdichtet die Erkenntnisse entlang des Untersuchungsgangs. Abschnitt~\ref{beantwortung} beantwortet die eingangs gestellte Forschungsfrage, benennt die Grenzen der Untersuchung und schlie{\ss}t die Arbeit ab.

\subsection{Zusammenfassung der zentralen Erkenntnisse}
\label{zusammenfassung}

Die Arbeit geht von einem Missverh\"altnis aus. Sprachverarbeitende Modelle erreichen eine Leistungsf\"ahigkeit, die sie f\"ur die textlastige Arbeit des First-Level-Supports geeignet erscheinen l\"asst, w\"ahrend die vorliegende Literatur \"uberwiegend einzelne Techniken oder abgegrenzte Anwendungsf\"alle betrachtet. Ein Konzept, das den gesamten Support-Prozess in den Blick nimmt und angibt, an welcher Stelle welche Aufgabe auf ein Modell \"ubergeht, fehlt weitgehend. Diese L\"ucke bestimmt den Zuschnitt der Arbeit.

Die Analyse des ITIL-basierten Referenzprozesses zeigt einen Ablauf, der zwischen Meldungseingang und Eskalationsentscheidung durchg\"angig auf menschlicher Arbeitsleistung ruht. Zwischen dem Eingang einer Meldung und dem ersten inhaltlichen L\"osungsversuch liegen mehrere Schritte, die den Vorgang lediglich einordnen, ohne zu seiner Behebung beizutragen. Als Kernschwachstellen treten zwei Felder hervor: die manuelle Kategorisierung, deren Treffsicherheit vom Erfahrungswissen des einzelnen Bearbeiters abh\"angt, sowie die unstrukturierten Nutzereingaben, die R\"uckfrageschleifen erzwingen und die stichwortbasierte Suche in den Wissensbest\"anden ins Leere laufen lassen. Beide Schwachstellen lassen sich an den etablierten Kennzahlen First Contact Resolution Rate, First Level Resolution Rate, Mean Time to Resolve und Ticketvolumen pro Agent verorten.

Das zentrale Ergebnis der Arbeit ist ein hybrider Soll-Prozess. Drei KI-gest\"utzte Schritte setzen an den identifizierten Schwachstellen an. Die Erfassung verl\"auft dialogorientiert und erkennt zugleich das Anliegen, die Kategorisierung und Priorisierung st\"utzt sich auf eine nach Aufgabentyp differenzierte Modellwahl, und die L\"osungsvorschl\"age entstehen abrufgest\"utzt aus dem vorhandenen Wissensbestand. Verankert sind diese Schritte in einem BPMN-Modell, das jeden Schritt einem maschinellen oder menschlichen Tr\"ager zuordnet und die Entscheidungsknoten samt konfidenz- und regelbasierter R\"uckfallkriterien ausweist. Datenschutz-Grundverordnung und KI-Verordnung wirken dabei als Design-Pr\"amissen, die den Zuschnitt dieser Knoten mitbestimmen, und nicht als nachgelagerte Pr\"ufpunkte.

Die W\"urdigung ergibt f\"ur die First Level Resolution Rate und die Mean Time to Resolve eine positive Wirkungsrichtung. F\"ur die First Contact Resolution Rate gilt dies nur prozessweit, denn im Agenten-Kanal kann sie sinken, wenn der Self-Service die einfachen F\"alle herausl\"ost; das Ticketvolumen pro Agent d\"urfte steigen, ist wegen des ver\"anderten Ticketmix aber nicht ohne Weiteres als Produktivit\"atsgewinn zu lesen. Zugleich begrenzen vier Felder den erreichbaren Automatisierungsgrad, n\"amlich die Neigung der Modelle zu nicht gedeckten Aussagen, die datenschutzrechtlichen Anforderungen an Ticketinhalte, die Angreifbarkeit \"uber den Freitext der Meldung und die Akzeptanz bei Agenten wie Anwendern. Aus ihnen folgt, dass die hybride Architektur keine Vorstufe der Vollautomatisierung darstellt, sondern deren begr\"undete Alternative.

\subsection{Beantwortung der Forschungsfrage und Schlusswort}
\label{beantwortung}

Die Untersuchung ging von folgender Frage aus:

\begin{quote}
\glqq Wie kann der First-Level-Support im Incident Management durch den Einsatz von Large Language Models prozessual automatisiert werden, sodass Effizienzpotenziale realisiert und zugleich Qualit\"ats-, Datenschutz- und Kontrollanforderungen gewahrt bleiben?\grqq{}
\end{quote}

Die drei Teilfragen aus Abschnitt~\ref{zielsetzung} beantwortet die Arbeit wie folgt.

\textbf{TF1.} Der Referenzprozess f\"uhrt von der Erfassung \"uber Kategorisierung und Priorisierung zur Erstdiagnose und endet entweder in einer L\"osung im First Level oder in der funktionalen Eskalation. Seine Schwachstellen liegen nicht in der L\"osungsarbeit, sondern in den ihr vorgelagerten Schritten. Die manuelle Kategorisierung h\"angt vom Erfahrungswissen des einzelnen Bearbeiters ab und erzeugt bei Fehlzuordnung R\"uckl\"aufer zwischen den Support-Ebenen; unstrukturierte Nutzereingaben erzwingen R\"uckfrageschleifen und lassen die stichwortbasierte Suche in den Wissensbest\"anden ins Leere laufen.

\textbf{TF2.} Der Soll-Prozess ist so zu gestalten, dass die Trennlinie zwischen Maschine und Mensch nicht entlang des technisch M\"oglichen verl\"auft, sondern entlang der Tragweite der jeweiligen Entscheidung. Maschinell erledigt werden reversible und risikoarme Schritte, deren Fehler im weiteren Verlauf auffallen und sich ohne bleibenden Schaden korrigieren lassen. Menschlich bleiben Entscheidungen, die den Vorgang beenden oder unmittelbar auf den Anwender wirken. Die Beteiligung des Menschen ist dabei nach der Verl\"asslichkeit des Vorschlags gestaffelt, und der Konfidenzschwellenwert bildet den Regler, \"uber den sich der Automatisierungsgrad mit belegter Vorschlagsqualit\"at erh\"ohen l\"asst.

\textbf{TF3.} Eindeutig positiv ist die erwartete Wirkungsrichtung f\"ur die First Level Resolution Rate und die Mean Time to Resolve. Die First Contact Resolution Rate ist nach der Umgestaltung nur getrennt nach Kan\"alen interpretierbar, weil der Self-Service die einfachen F\"alle aus dem Agenten-Kanal herausl\"ost, und ein steigendes Ticketvolumen pro Agent spiegelt nicht ohne Weiteres einen Produktivit\"atsgewinn wider. Die Grenzen liegen in den vier in Abschnitt~\ref{zusammenfassung} genannten Feldern, die den Automatisierungsgrad aus technischen, rechtlichen und organisatorischen Gr\"unden begrenzen.

Auf die Hauptfrage folgt daraus die nachstehende Antwort. Das Konzept sieht die Automatisierung dort vor, wo die der eigentlichen L\"osung vorgelagerten Schritte an KI-Komponenten \"ubergehen. Erfassung, Anliegenserkennung, Kategorisierung und Priorisierung verlangen kein Fachurteil \"uber die St\"orung, sondern die Verarbeitung von Text, und genau darin liegt die St\"arke der Modelle. Der Agent tritt dadurch nicht fr\"uher, sondern sp\"ater in den Prozess ein, und zwar dort, wo eine Entscheidung zu treffen ist. Die erwartete Effizienzwirkung ergibt sich folglich nicht daraus, dass ein Modell die Arbeit des Menschen nachahmt, sondern daraus, dass es ihm die Vorarbeit abnimmt.

Die Qualit\"ats-, Datenschutz- und Kontrollanforderungen werden im Konzept nicht \"uber eine nachgelagerte Pr\"ufung adressiert, sondern im Prozess selbst verankert. Die Beteiligung des Menschen an definierten Knoten, die Konfidenzschwelle vor der automatischen \"Ubernahme eines Vorschlags, die Positivliste f\"ur vollautomatisiert beantwortbare F\"alle und der dokumentierte Kontexttransfer bei der Eskalation sind Bestandteile des Modells und keine begleitenden Ma{\ss}nahmen. Darin liegt die zentrale Aussage der Arbeit. Die Prozessmodellierung ist hier nicht Dokumentation eines ohnehin feststehenden Ablaufs, sondern Gestaltungsinstrument; erst sie macht sichtbar und damit verhandelbar, welche Entscheidung eine Maschine und welche ein Mensch trifft.

Die Reichweite dieser Antwort ist aus den in Abschnitt~\ref{zielsetzung} genannten Gr\"unden begrenzt. Hinzu tritt eine Einschr\"ankung, die in der gew\"ahlten Vorgehensweise selbst liegt, denn Konstrukteur und Bewerter des Artefakts sind dieselbe Person, sodass die Evaluation argumentativ bleibt und keine Au{\ss}ensicht einbezieht. Gemessen an den Kriterien des Memorandums zur gestaltungsorientierten Wirtschaftsinformatik liegt der Beitrag der Arbeit daher weniger in einem belegten Nutzen als in der Abstraktion, n\"amlich einem Referenzkonzept, das sich auf unterschiedliche Support-Organisationen \"ubertragen l\"asst.

Daraus ergeben sich drei Anschlussm\"oglichkeiten f\"ur weitere Forschung. Eine Pilotierung mit Vorher-Nachher-Messung der vier Kennzahlen w\"urde die hier abgeleiteten Wirkungsrichtungen empirisch pr\"ufen. Eine Validierung des Referenzprozesses durch Fachleute aus dem Service-Desk-Betrieb w\"urde die fehlende Au{\ss}ensicht erg\"anzen. Eine \"Ubertragung des Konzeptes auf das Service-Request-Management, das die vorliegende Arbeit ausklammert, w\"urde seinen Geltungsbereich erweitern.

Der entwickelte Rahmen bleibt f\"ur eine weitergehende Automatisierung anschlussf\"ahig, da sich wachsende Autonomie \"uber die vorhandenen Schwellen und Kataloge steuern l\"asst, ohne die Prozessstruktur zu ver\"andern. Die Gestaltungsaufgabe verschiebt sich damit dauerhaft. Sie besteht weniger in der Entscheidung, ob im First-Level-Support automatisiert wird, als in der Festlegung, an welchen Stellen die Entscheidungshoheit des Menschen verankert bleibt.
